\section{习近平法治思想}
\begin{enumerate}
  \item \zd{政治题}。看个 20 多遍。800 - 1000 字。\bs{不要套模板}。15 年开始做。
        \begin{enumerate}
          \item 单选。\zd{很少负面选项}。主要选错的。多选。主要选对的。
          \item \bs{不能用正确压倒精确、用现实去否定理论}。
        \end{enumerate}
  \item 理论背景。\zd{全面}依法治国。
  \item 形成发展。时代背景、\zd{3 大逻辑}、历史进程、9 大节点、\zd{5 大特色}。
        \begin{enumerate}
          \item 2020 年 11 月，明确\zd{指导地位}。
          \item \zd{制度威力，科学指南}。
          \item 3 大逻辑。理论、历史、实践。
          \item 5 个特色。\bs{时代、人民、实践、原创、系统}。时代性、人民性、实践性、原创性、系统性。
        \end{enumerate}
  \item 概念、地位。背。
        \begin{enumerate}
          \item 着眼于中华民族伟大复兴战略全局和当今世界百年未有之大变局。
          \item 顺应中华民族伟大复兴时代要求应运而生的重大战略指导思想。
          \item 习近平新时代中国特色社会主义思想的重要组成部分。
        \end{enumerate}
  \item \bs{重大意义}。主观题。4 个关键句。
        \begin{enumerate}
          \item 是马克思主义法治理论同中国实际相结合的\zd{最新成果}。
          \item 是对党领导法治建设丰富实践和宝贵经验的\zd{科学总结}。
          \item 是引领新时代法治建设的\zd{思想旗帜}。
          \item 是全面依法治国的\zd{根本遵循}。
        \end{enumerate}
  \item 核心要义。俗称 \zd{11 个大坚持}。
        \begin{enumerate}
          \item \zd{党}对全面依法治国的领导。
          \item 以\zd{人}民为中心。
          \item 中国\zd{特色}社会主义法治道路。
          \item \zd{两依宪}。依宪治国、依宪执政。
          \item \zd{立、执、司、守没有监}。全面推进\zd{科学立法、严格执法、公正司法、全民守法}。
          \item \zd{三依三治}。依法治国、依法执政、依法行政共同推进。法治国家、法治政府、法治社会一体建设。
          \item 在法治轨道上推进国家治理体系和治理能力\zd{现代化}。
          \item 抓住领导\zd{干部}这个“关键少数”。
          \item 建设德才兼备的高素质法治工作\zd{队伍}。
          \item 统筹推进国\zd{内}法治和涉\zd{外}法治。
          \item \zd{抓}。建设中国特色社会主义法治体系。
        \end{enumerate}
  \item 实践要求。5 大作用，4 对关系。
        \begin{enumerate}
          \item 五位一体 + 保障。法治保障\zd{政治}稳定、\zd{经济}发展、\zd{文化}繁荣、\zd{社会}和谐、\zd{生态}良好。
          \item 法治和政治、法治和改革、依法治国和以德治国、依法治国和依规治党。
        \end{enumerate}
\end{enumerate}
